\subsection{A.2(a): a cap cannot depend only on its own length}

\paragraph{Construction and source comparison.}
Here a cap is the embedded pointed disc obtained by attaching a flat
semi-infinite cylinder along a geodesic boundary by arclength, then filling
its puncture. Its boundary is the original geodesic seam. This is the
infinite-grafting construction in Mondello (2011), Theorem~5.4
(the downloaded source has the label \texttt{thm:grafting}), restricted
to the smooth finite stratum. F\i rat (2021), equations~(1.2), (1.3a),
and the projective-connection and cap-coordinate formulas in \S3,
describes the same construction
on a pair of pants: the boundary lengths are $L_i=2\pi\lambda_i$,
the cap coordinates are $w_i$, and the flat metric on a cap is
$\lambda_i^2|dw_i|^2/|w_i|^2$. We opened both sources and their
downloaded \LaTeX{} files. We now give a direct argument that does not
require choosing branches of the gamma functions in the latter paper.

\begin{theorem}[Proved: a single cap remembers all the lengths]
\label{A:single-cap-all-lengths}
Let $(X;p_1,\ldots,p_n)$ be a fixed connected compact marked Riemann
surface. Suppose two systems of pairwise disjoint embedded pointed
discs $(D_j,p_j)$ and $(D'_j,p_j)$ arise by infinite grafting of
hyperbolic surfaces with geodesic boundaries of positive lengths
$\ell_j$ and $\ell'_j$, respectively. If, for one index $i$,
\[
 D_i=D'_i,
\]
then $\ell_j=\ell'_j$ for every $j$, including $j=i$.
With the usual marking convention,
Mondello's grafting theorem further implies that the two complete
systems are equal.
\end{theorem}

\begin{proof}
A hyperbolic surface carries its Fuchsian complex projective structure:
on its universal cover a local isometry into the upper half-plane is a
holomorphic developing map and two choices differ by an element of
$\operatorname{PSL}_2(\mathbb R)$. Near a lift of a geodesic boundary
of length $\ell$, choose a logarithmic developing coordinate
$\zeta=t+is$ such that the boundary is $s=0$, its arclength coordinate
is $t$, its deck transformation is $\zeta\mapsto\zeta+\ell$, and
the developing map is $i e^{\zeta}$ up to a M\"obius transformation.
The direction of $s$ can be chosen to point into the attached cylinder;
reversing this convention only reverses the sign of the exponent below.
The grafted cylinder has exactly the same complex coordinate $\zeta$.
Consequently $i e^{\zeta}$ extends the projective structure across the
seam and throughout the cylinder. It is locally univalent everywhere
on the cylinder. Its monodromy under one turn is multiplication by
$e^{\ell}$, a projective transformation, so this continuation defines
a projective structure on the punctured capped surface.

On the filled cap let
\[
 w=\exp(2\pi i\zeta/\ell),\qquad
 \lambda=\ell/(2\pi).
\]
Thus $w$ maps the cap to the unit disc, takes its centre to $0$, and
is determined by the pointed embedded disc up to multiplication by a
constant of absolute value~1. The developing map on any simply
connected subregion of the punctured cap is, up to a projective change
of coordinate, $w^{-i\lambda}$. A different sign convention gives
$w^{i\lambda}$ and the same Schwarzian derivative.

Write $\{f,z\}=f'''/f'-\frac32(f''/f')^2$. Differentiating
$w^\alpha$ gives
\[
 \{w^\alpha,w\}=\frac{1-\alpha^2}{2w^2}.
\]
The Schwarzian chain rule therefore gives the projective connection
in any holomorphic coordinate $z$ on the cap as
\begin{equation}\label{A:cap-projective-connection}
 Q(z)=\{w,z\}+
 \frac{1+\lambda^2}{2}\left(\frac{w'(z)}{w(z)}\right)^2.
\end{equation}
This expression is invariant under the allowed rotations of $w$.
Since $w$ has a simple zero at the centre and $w'$ does not vanish,
$\{w,z\}$ is holomorphic there, whereas the second term has a pole
of order two with leading coefficient $(1+\lambda^2)/2$. We have
thus proved, rather than assumed, that the connection extends to a
meromorphic projective connection on $X$ with regular singularities
at its marked points and with precisely these leading coefficients.
There are no other poles: the developing maps on the hyperbolic core,
the collars and the cylinders are locally univalent.

Apply this construction to the two cap systems. We first prove the
stronger assertion that their common cap determines its own length.
Suppose $D_i=D'_i$ but $\lambda_i\ne\lambda'_i$, and use their
common normalized coordinate $w$, unique up to rotation. Since
$Q-Q'$ is a globally defined meromorphic quadratic differential,
so is
\[
 \Phi=\frac{Q-Q'}{(\lambda_i^2-(\lambda'_i)^2)/2}.
\]
On the cap, $\Phi=(dw/w)^2$. Choose a developing map $u$ of the
first projective structure, with $u=w^{i\lambda_i}$ on a simply
connected portion of this cap. Then
\[
 (du/u)^2=-\lambda_i^2\Phi
\]
there, hence on the whole universal cover of the punctured surface
by meromorphic continuation. If $M$ is any projective monodromy
transformation, deck invariance of the right side implies
\[
 \left(\frac{d(M(u))}{M(u)}\right)^2=(du/u)^2.
\]
The developing image contains an open set, so the rational function
identity $(M'(z)/M(z))^2=z^{-2}$ follows. On a connected open set
its square root has a constant sign, giving
$M'(z)/M(z)=1/z$ or $-1/z$. Integration, and then rational
continuation, gives respectively $M(z)=az$ or $M(z)=a/z$.
Thus all holonomy lies in the normalizer of $\{0,\infty\}$,
which has an abelian subgroup of index two.

This is impossible. Attaching punctured cylinders is a homotopy
equivalence, so this holonomy contains the faithful Fuchsian
holonomy of the hyperbolic core. A finite-area hyperbolic core with
geodesic boundary has $\chi<0$ by Gauss--Bonnet. With a nonempty
boundary its fundamental group is free of rank
$2g+n-1\ge2$. Such a group is not virtually abelian: an index-$d$
subgroup is the fundamental group of a $d$-sheeted covering of a
rank-$r$ graph, and has rank $1+d(r-1)\ge2$. This contradiction
proves $\lambda_i=\lambda'_i$.

On the common cap $D_i$, the two normalized disc coordinates differ
by a rotation.
Equation~\eqref{A:cap-projective-connection} and
$\ell_i=\ell'_i$ show that the two projective connections coincide
there. Their difference is a globally defined meromorphic quadratic
differential on $X$: the Schwarzian coordinate-change terms cancel.
It vanishes on a nonempty open subset, hence vanishes everywhere by
the identity theorem. Comparing its leading double-pole coefficients
at each $p_j$ gives
\[
 1+(\ell_j/2\pi)^2=1+(\ell'_j/2\pi)^2.
\]
All lengths are positive, so $\ell_j=\ell'_j$ for every $j$.
Finally the restriction of Mondello's map
$(\mathrm{gr}_\infty,L)$ to positive finite lengths is injective;
the fixed capped marked surface and the common length vector therefore
determine the same marked bordered surface and its attached caps.
This last invocation is needed only for equality of the whole systems,
not for the equality of their lengths just proved.
\end{proof}

\begin{corollary}[Proved extension: positive caps distinguish cuspidal ends]
\label{A:single-cap-with-cusps}
The preceding uniqueness holds also when some of the other ends are
complete hyperbolic cusps, assigned length zero, while the compared
cap is a positive cap in both constructions. Therefore keeping one
positive pointed cap unchanged cannot be compatible with replacing
any other positive geodesic boundary by a complete cusp on the same
marked compactification.
\end{corollary}

\begin{proof}
At a complete cusp a parabolic developing coordinate is, up to a
projective transformation and a nonzero scalar, $\log w$, where
$w$ is a holomorphic cusp coordinate vanishing simply at the filled
point. Direct differentiation gives
$\{\log w,w\}=1/(2w^2)$. Hence the projective connection extends
meromorphically there with the same leading coefficient
$(1+\lambda^2)/2$ at $\lambda=0$. The compared cap is positive,
so the holonomy argument above applies unchanged and first recovers
its own positive length. The identity theorem then identifies the
two global connections. Their double-pole coefficients recover all
the nonnegative lengths, including which ones are zero. In
particular a positive length cannot become zero.
\end{proof}

\begin{proposition}[Refuted under the literal cusp admissibility convention]
\label{A:MDB-forgetting-admissibility}
Interpret the admissibility of the external cap system in
67/075(b) to mean that on each normalization component the node
branches are complete cusps, as the point-like boundary components
$s,t$ on 67/081 specify. Under this convention, the proposed
capping-and-forgetting operation
$MDB_{g,n}\to M_{\mathrm{st}}ND_{g,n}$ need not take values in
its stated target. It already fails for two positive-boundary
pairs of pants glued along one boundary, giving type $(0,4)$.
More generally it fails on any normalization component that has
both a forgotten positive node cap and a retained positive
external cap.
\end{proposition}

\begin{proof}
Choose two hyperbolic pairs of pants with three positive boundary
lengths each, the two boundaries to be glued having equal lengths.
Glue these geodesic boundaries by an arclength-reversing isometry.
This produces a compact bordered hyperbolic four-holed sphere with
a specified separating geodesic; hence it is a legitimate source
object. Cut it again and attach its canonical infinite cylinders.
The compactified components are three-pointed spheres, each
carrying a full system of three positive canonical caps. Forget
the cap at the branch corresponding to the separating geodesic,
but keep the two external caps and the node's nonzero tangent
tensor, exactly as prescribed on 67/075.

If the remaining caps were admissible with the node branch treated
as a cusp, the same compactified three-pointed sphere would be
obtained by grafting a second hyperbolic surface with that end of
length zero and with the retained external caps unchanged.
Choose either retained positive cap. Corollary~\ref{A:single-cap-with-cusps}
forces equality of all three lengths in the two constructions,
contradicting positivity of the first construction's node-boundary
length and zero length of the second. The same argument is local
to a normalization component and proves the general assertion.
\end{proof}

\paragraph{The two meanings of a partial admissible system.}
This is a conditional interpretation result, not an assertion that
the manuscript formally defines its target category. The wording
on 67/075(b) requires an admissible system only around external
marks, and the 67/081 pants target explicitly changes the glued
boundaries into point-like components. We checked the latter phrase
against the original scan. A different, precise definition is to
call a partial cap system \emph{extendible} when it occurs by
forgetting selected caps from a full admissible system. That
definition makes the displayed forgetting operation well-defined,
but it does not identify the retained caps with the canonical caps
obtained after declaring omitted ends to be cusps. With the nodal
tensor still independent and unrestricted, the radius obstruction
of Theorem~\ref{A:04} remains: even an extendible external system
cannot produce a node tensor of norm $17$ in its fixed normalized
three-pointed coordinates, since every possible node cap has
conformal radius at most $4$ and their product at most $16$.

\begin{corollary}[Refuted: the question marked ``?'' at 67/068]
\label{A:pants-cap-dependence}
Normalize the capped, labelled pair of pants to
$(\mathbb P^1;0,1,\infty)$. The cap at $0$ is not a function of
its own boundary length alone. Specifically the length vectors
\[
 (2\pi,2\pi,2\pi),\qquad(2\pi,4\pi,2\pi)
\]
give different embedded pointed discs at $0$, although that boundary
has the same length in both examples. Every positive length vector is
realized by a hyperbolic pair of pants, so both examples meet all
hypotheses of the question.
\end{corollary}

\begin{proof}
The existence statement follows either from the usual gluing of two
right-angled hyperbolic hexagons with alternating side lengths
$\ell_1/2,\ell_2/2,\ell_3/2$, or directly from the positive-length
stratum of the opened grafting theorem. The three filled punctures
have the unique labelled normalization $0,1,\infty$. If the two caps
at $0$ were the same, Theorem~\ref{A:single-cap-all-lengths} would
force the second boundary lengths to be equal, a contradiction.

For an explicit local witness, put
$\Delta_j=(1+\lambda_j^2)/2$. A projective connection on the
three-punctured sphere with only these regular singularities must be
\begin{equation}\label{A:pants-projective-connection}
 Q(z)=\frac{\Delta_1}{z^2}
      +\frac{\Delta_2}{(z-1)^2}
      +\frac{\Delta_3-\Delta_1-\Delta_2}{z(z-1)}.
\end{equation}
Indeed the difference of two such connections has at most simple
poles at the three marked points, and would be a section of
$K_{\mathbb P^1}^{\otimes2}(0+1+\infty)$, of degree $-1$;
there is no nonzero such section. Alternatively partial fractions
and the coordinate $1/z$ verify the displayed formula directly.
Let $w(z)=c(z+a z^2+O(z^3))$ be the unit-disc coordinate on the
cap at $0$, where $c\ne0$. Its Schwarzian is holomorphic at $0$,
and $(w'/w)^2=z^{-2}+2a z^{-1}+O(1)$.
Comparison of the simple-pole coefficients in
\eqref{A:cap-projective-connection} and
\eqref{A:pants-projective-connection} yields
\[
 a=\frac{\Delta_1+\Delta_2-\Delta_3}{2\Delta_1}
  =\frac{1+\lambda_1^2+\lambda_2^2-\lambda_3^2}
         {2(1+\lambda_1^2)}.
\]
The two examples give $a=1/2$ and $a=5/4$. The ratio of the
quadratic and linear coefficients is unchanged by rotating $w$;
equal pointed discs would have equal ratios. This is also exactly
the independently opened expansion in F\i rat (2021), equation
(1.3a). The program \texttt{code/A567\_matrices.py} computes these
two rational numbers exactly, and its JSON output records the result.
\end{proof}

\paragraph{Scope.}
This proves the stated dependence on the other lengths; a single
positive cap in fact determines the entire length vector. It does
not by itself establish that
the vector of conformal radii determines the cap system: a conformal
radius discards the higher Taylor coefficients used in this proof.
Nor does the argument identify a one-variable limiting cap when several
lengths vary without a specified direction.
